\todo{Lecture 13}

\subsection{Fonctions lipschitziennes}
\begin{definition}
On dit que $f:D\to \mathbf{R}$ est $L$-lipschitzienne pour $L\ge 0$ si 
$$
\forall x,y\in D:|f(x)-f(y)|\le L|x-y|.
$$
On appelle $L$ une constante de Lipschitz pour $f$. On dit que $f$ est lipschitzienne s'il existe un tel $L$.
\end{definition}
\begin{remark}
Etre lipschitzienne implique uniformement continue implique continue. Mais pas dans le sens inverse !
\end{remark}
\begin{definition}
$f$ est contractante si elle est $L$-lipschitzienne avec $L<1$. 
\end{definition}
\begin{theorem}[Point fixe de Banach]
Si $f:\mathbf{R}\to \mathbf{R}$ est contractante, alors il existe un unique $x \in \mathbf{R}$ tel que $f(x)=x$.
\end{theorem}
\begin{proof}
\begin{itemize}
\item Unicite : Soient $u,v\in \mathbf{R}$ tels que $f(u)=u$ et $f(v)=v$. Alors
$$
|u-v|=|f(u)-f(v)|\le L|u-v|.
$$
Donc $(1-L)|u-v|\le 0$ mais $(1-L)>0$ donc $|u-v|=0$.
\item Existence : Fixons $x_{0}\in \mathbf{R}$. On definie $x_{n+1}=f(x_{n})$ pour $n\in \mathbf{N}$. Regardons la distance entre les points de $(x_{n})$. 
\begin{align*}
|x_{2}-x_{1}| & =|f(x_{1})-f(x_{0})|\le L|x_{1}-x_{0}| \\
|x_{3}-x_{2}| & =|f(x_{2})-f(x_{1})|\le L|x_{2}-x_{1}|\le L^{2}|x_{1}-x_{0}| \\
 & ~\vdots \\
 |x_{n+1}-x_{n}| &\le L^{n}|x_{1}-x_{0}| 
\end{align*}
donc 
\begin{align*}
|x_{n+m}-x_{n} | & \le |x_{n+m}-x_{n+m-1}|+|x_{n+m-1}-x_{n+m-2}|+\dots +|x_{n+1}-x_{n}| \\
&\le(L^{n+m-1}+L^{n+m-2}+\dots +L^{n})|x_{1}-x_{0}| \\
&\le \frac{L^{n}}{1-L}|x_{1}-x_{0}|.
\end{align*}
Donc $(x_{n})$ est de Cauchy et converge donc vers un $\overline{x}\in \mathbf{R}$. Donc
$$
\overline{x}=\lim_{ n \to \infty } x_{n}=\lim_{ n \to \infty }x_{n+1}=\lim_{ n \to \infty } f(x_{n})=f(\overline{x}).
$$
\end{itemize}
\end{proof}
\chapter{Suites de fonctions}
\begin{definition}
Soit $f:D\to \mathbf{R}$ et $(f_{n})$ une suite de fonctions sur $D$. On dit que $(f_{n})$ converge simplement (ou ponctuellement) vers $f$ si 
$$
\forall x \in D:\lim_{ n \to \infty } f_{n}(x)=f(x).
$$
Pour chaque $x$ fixe, il s'agit de la limite d'une suite reelle.
\end{definition}
\begin{notation}
Si $(f_{n})$ converge simplement vers $f$, on note
$$
f_{n}\overset{\mathrm{ponct}}\longrightarrow f.
$$
\end{notation}
\begin{definition}
On dit que $(f_{n})$ converge uniformement vers $f$ si
$$
\forall\varepsilon>0\,\exists N\in \mathbf{N}\,\forall n\ge N,\,\forall x \in D:|f_{n}(x)-f(x)|<\varepsilon.
$$
\end{definition}
\begin{notation}
Si $(f_{n})$ converge uniformement vers $f$, on ecrit
$$
f_{n}\overset{\mathrm{unif}}\to f.
$$
\end{notation}
\begin{theorem}
Soient $f:D\to \mathbf{R}$ et $(f_{n})$ une suite dans $C^{0}(D)$. Si $(f_{n})$ converge uniformememnt vers $f$, alors $f\in C^{0}(D)$.
\end{theorem}
\begin{proof}
Fixons $x_{0}\in D$ et $\varepsilon>0$. Par convergence uniforme,
$$
\exists N\,\forall n\ge N\,\forall x \in D:|f_{n}(x)-f(x)|<\varepsilon /3.
$$
Donc
\begin{align*}
|f(x)-f(x_{0})| & =|f(x)-f_{n}(x)+f_{n}(x)-f_{n}(x_{0})+f_{n}(x_{0})-f(x_{0})| \\
 & \le \underbrace{ |f(x)-f_{n}(x)| }_{ <\varepsilon /3 }+|f_{n}(x)-f_{n}(x_{0})|+\underbrace{ |f_{n}(x_{0})-f(x_{0})| }_{ <\varepsilon /3 }.
\end{align*}
Comme $f_{n}$ est continue, alors il existe $\delta>0$ tel que $|x-x_{0}|<\delta$ et donc $|f_{n}(x)-f_{n}(x_{0})|<\varepsilon /3$.
En combinant, on a donc montrer que pour tout $x_{0}\in D$ et pour tout $\varepsilon>0$ on a que $f$ et continue sur $D$.
\end{proof}